% Chapter 20: Perfectoid Spaces and the Revolution in p-adic Geometry
\let\LocalMacro\relax

\chapter{Perfectoid Spaces and the Revolution in p-adic Geometry}

\section{The Problem}

\subsection{Informal}

\textbf{Background.} Number theory studies equations over two fundamentally different kinds of number systems: those of characteristic 0 (like the ordinary integers and rational numbers) and those of characteristic $p$ (like clock arithmetic with a prime number of hours). These two worlds have different tools and different structures, and for decades there was no way to translate between them.

\textbf{Given.} You are given a geometric object (a space of solutions to equations) defined over a $p$-adic number field---a number system that extends the integers using a prime $p$.

\textbf{Goal.} Build a universal translation tool that converts problems from one number system (the $p$-adics) into equivalent problems in a completely different number system (characteristic $p$), so that hard tools from one side can be applied to the other.

\textbf{Constraints.} The bridge must work not just for individual spaces but for families of spaces, including singular and non-projective ones. It must preserve the cohomological invariants that number theorists care about. The construction must be functorial---it must respect the relationships between spaces.

\textbf{Example.} For a single smooth projective variety over a $p$-adic field, Fontaine's $p$-adic Hodge theory already provided a dictionary to characteristic $p$. The breakthrough was extending this to \emph{families} of spaces, which required the entirely new framework of perfectoid spaces.

\subsection{Formal}

The central question is: can one construct a geometric framework that provides a ``dictionary'' between $p$-adic geometry (mixed characteristic $(0,p)$) and geometry over $\mathbb{F}_p((t))$ (equal characteristic $p$)? Specifically, can one build a category of spaces over a $p$-adic field that is equivalent to a corresponding category over characteristic $p$, preserving cohomological information?

\subsection{Historical Context}
Peter Scholze, then a 23-year-old PhD student at Bonn, invented a new kind of geometric object---perfectoid spaces---that serves as a universal bridge between these two languages. The work earned him the 2018 Fields Medal and has reshaped modern number theory.

\subsection{What Was Known}
Before Scholze, mathematicians had partial translation tools between mixed characteristic and equal characteristic $p$ geometry, but they were limited to special cases. Fontaine's theory of p-adic Hodge theory worked for individual spaces but not for families. The pro-\'etale topology, which Scholze later developed, did not yet exist. The full power of perfectoid spaces was completely new.

\subsection{Why It Matters}
Perfectoid spaces provide a universal bridge between two fundamentally different ways of studying equations. This bridge has already led to major breakthroughs in number theory, including progress on the p-adic Langlands program and the weight-monodromy conjecture. The techniques have also spawned new fields like prismatic cohomology and condensed mathematics.

\subsection{Simplified Version}
For individual spaces, the translation between characteristic 0 and characteristic $p$ was already understood: Fontaine's $p$-adic Hodge theory provided a dictionary for smooth projective varieties over $p$-adic fields. The breakthrough was making this dictionary work for \emph{families} of singular, non-projective spaces---which required an entirely new geometric framework.

\section{Mathematical Theory}


\subsection{Prerequisites}
This chapter assumes familiarity with:
\begin{itemize}
\item Algebraic number theory (local fields, ramification)
\item Rigid analytic geometry (Tate algebras, affinoid spaces)
\item \'Etale cohomology and Galois representations
\item Basic category theory (limits, adjoint functors)
\end{itemize}

\subsection{Perfectoid Fields}
The key insight: many problems in p-adic geometry become trivial in the limit of \emph{perfect} fields.

\begin{definition}[Perfectoid Field]
A complete non-archimedean field $K$ of characteristic $0$ with residue characteristic $p>0$ is \emph{perfectoid} if:
\begin{enumerate}
\item The value group $|K^\times|$ is dense in $\mathbb{R}_{>0}$.
\item The Frobenius map $\mathcal{O}_K/p \to \mathcal{O}_K/p$ is surjective.
\end{enumerate}
\end{definition}

\begin{exercise}[Basic]
Show that $\mathbb{C}_p$ is perfectoid. Why is $\mathbb{Q}_p$ not perfectoid?
\end{exercise}

\subsection{The Tilt Equivalence}
Scholze's fundamental construction: to any perfectoid field $K$ of characteristic $0$, associate a perfectoid field $K^\flat$ of characteristic $p$.

\begin{definition}[Tilt]
$K^\flat = \varprojlim_{x \mapsto x^p} K$, with valuation $|(x^{(n)})| = |x^{(0)}|$.
\end{definition}

\begin{theorem}[Fontaine-Wintenberger / Scholze]
There is an equivalence of categories:
\[
\{\text{finite extensions of } K\} \simeq \{\text{finite extensions of } K^\flat\}
\]
preserving Galois groups and ramification.
\end{theorem}

This is the \emph{tilt equivalence}: it transfers problems from mixed characteristic $(0,p)$ to equal characteristic $p$, where powerful tools (Frobenius, Cartier operator) are available.

\subsection{Perfectoid Spaces}
\begin{definition}[Perfectoid Space]
A locally ringed space $(X, \mathcal{O}_X)$ that is locally isomorphic to $\text{Spa}(R, R^+)$ for a perfectoid Tate-Huber pair $(R, R^+)$.
\end{definition}

Key properties:
\begin{itemize}
\item The category of perfectoid spaces contains all rigid analytic spaces over $\mathbb{Q}_p$ and $\mathbb{F}_p((t))$.
\item There is a \emph{tilt functor} $(-)^{\flat}$ from char $0$ to char $p$ perfectoid spaces.
\item \emph{Pro-\'etale site}: Scholze constructed a Grothendieck topology where $\mathcal{O}_X^+$ is acyclic, enabling p-adic Hodge theory in families.
\end{itemize}

\section{The Solution}

\subsection{The Big Idea}
Number theory has two kinds of geometry: one over $p$-adic fields (characteristic 0, where $p$ is a prime) and one over $\mathbb{F}_p$ (characteristic $p$, where $p = 0$). Fontaine's $p$-adic Hodge theory provided a dictionary between them for individual smooth varieties, but the dictionary broke down for \emph{families} of spaces. Scholze's insight was to find the right algebraic operation---the \emph{tilt}---that takes any perfectoid field $K$ of characteristic 0 and produces a field $K^\flat = \varprojlim_{x \mapsto x^p} K$ of characteristic $p$ carrying exactly the same Galois groups and ramification data. The surprise: this is not just a field-level trick. Scholze showed it extends to \emph{spaces}: every perfectoid space of characteristic 0 has a tilt that is a perfectoid space of characteristic $p$, and the two have the same étale cohomology. The three ingredients are: (1)~the \textbf{tilt equivalence} $K \mapsto K^\flat$ at the field level, (2)~the category of \textbf{perfectoid spaces} that carry the tilt functor globally, and (3)~the \textbf{pro-étale topology} that makes $\mathcal{O}_X^+$ acyclic, so comparison isomorphisms extend to families. The result: a geometric framework where mixed characteristic and equal characteristic $p$ geometry are genuinely the same theory.

\subsection{What Was Stopping Progress}
The existing dictionary (Fontaine's $p$-adic Hodge theory) could translate individual sentences---a single smooth variety over a $p$-adic field could be compared with its characteristic-$p$ shadow. But when you tried to translate \emph{families} of spaces---saying something about an entire family at once---the dictionary broke down. It was like having a word-by-word translator who can handle individual sentences but freezes when confronted with an entire novel. The tools of rigid analytic geometry were too rigid (pun intended): their underlying ``topology'' was too coarse to detect the fine arithmetic phenomena that needed translating.

\subsection{The Breakthrough}
Scholze introduced three interlocking ideas that together replace the old toolkit entirely. First, the \textbf{tilt equivalence}: a construction that takes any perfectoid field of characteristic 0 and produces a perfectoid field of characteristic $p$ carrying exactly the same arithmetic information---like a photograph that captures every detail of its subject. Second, the category of \textbf{perfectoid spaces}: geometric objects that exist natively in both characteristics. Third, the \textbf{pro-étale topology}: a new way of ``gluing'' these spaces together that makes the dictionary work globally, for families of spaces, not just individual ones. The result: a Rosetta Stone for number theory.

\subsection{How It Works}

\subsection{What Was Missing}
For decades, mathematicians had a ``dictionary'' between characteristic~$0$ and characteristic~$p$ geometry, but only for individual, well-behaved spaces. Fontaine's $p$-adic Hodge theory gave a precise comparison between de Rham and crystalline cohomology for a \emph{single} smooth proper variety over a $p$-adic field. What was missing was any framework that could handle \emph{families}---families of singular, non-projective spaces whose geometry changes from fibre to fibre. Without such a framework, fundamental questions (like the weight-monodromy conjecture) could not even be stated globally, let alone proved.

\subsection{Why Existing Methods Failed}
The classical tools---rigid analytic geometry and the étale topology---were designed for individual spaces. The étale site of a rigid space is too coarse: it cannot detect $p$-adic Hodge-theoretic phenomena, and $\mathcal{O}_X^+$ has non-trivial cohomology. Attempts to extend Fontaine's constructions to families ran into walls: comparison isomorphisms would break down, filtrations would become ill-defined, and there was no cohomology theory that was simultaneously $p$-adic and well-behaved for non-smooth spaces. In short, the existing geometric foundations were not powerful enough.

\subsection{What Was Novel}
Scholze's innovation was threefold: (1)~the \textbf{tilt equivalence}, which associates to every perfectoid field of characteristic~$0$ a perfectoid field of characteristic~$p$ while preserving all arithmetic data; (2)~the category of \textbf{perfectoid spaces}, which are locally isomorphic to adic spaces built from perfectoid rings and carry a natural tilt functor between characteristic~$0$ and characteristic~$p$; and (3)~the \textbf{pro-étale topology}, a new Grothendieck topology in which $\mathcal{O}_X^+$ becomes acyclic and comparison isomorphisms extend to families. Together these three innovations replace the rigid-analytic toolkit with a far more flexible foundation.

\subsection{Student-Friendly Explanation}
Imagine you have two languages---one for describing geometry over numbers like $0, 1, 2, \ldots$ (characteristic~$0$) and one for geometry over numbers like $0, 1, 1+1, \ldots$ where eventually $p=0$ (characteristic~$p$). For a long time, people had a pocket dictionary that worked for individual ``words'' (individual spaces), but it broke down when you tried to translate whole ``paragraphs'' (families of spaces). Scholze built a universal translator by first introducing a perfectoid space, which lives in both languages simultaneously, and then proving that the two languages are secretly the same---you just need the right topology to see it. This is like discovering that two apparently unrelated writing systems are actually encoding the same story.

\subsection{Scholze's Thesis (2011) and Annals Papers (2012-2013)}

\begin{theorem}[Scholze \cite{scholze2012} -- Local Langlands for $\text{GL}_n$]
The local Langlands correspondence for $\text{GL}_n$ over a p-adic field can be realized geometrically in the \'etale cohomology of \emph{Lubin-Tate perfectoid spaces}.
\end{theorem}

\begin{theorem}[Scholze \cite{scholze2013} -- p-adic Hodge Theory for Rigid Spaces]
For a smooth proper rigid space $X$ over a p-adic field, there is a natural comparison isomorphism:
\[
H^i_{\text{\'et}}(X_{\mathbb{C}_p}, \mathbb{Z}_p) \otimes_{\mathbb{Z}_p} \mathbf{B}_{\text{dR}} \cong H^i_{\text{dR}}(X) \otimes_K \mathbf{B}_{\text{dR}}
\]
compatible with filtrations and Galois actions. This extends Fontaine's theory to families.
\end{theorem}

\begin{highlightbox}
The perfectoid approach turned the \emph{weight-monodromy conjecture} (a major open problem in p-adic Hodge theory) into a geometric statement about the \'etale cohomology of perfectoid spaces, which was subsequently proved by Bhatt, Morrow, and Scholze (2016) using the \emph{pro-\'etale topology}.
\end{highlightbox}

\subsection{Subsequent Developments (2014--2025)}
\begin{itemize}
\item \textbf{Bhatt-Morrow-Scholze (2016)}: Integral p-adic Hodge theory via $A_{\text{inf}}$-cohomology.
\item \textbf{Scholze-Weinstein (2016)}: Moduli of p-divisible groups and the Rapoport-Zink space as a perfectoid space.
\item \textbf{Fargues-Fontaine curve (2018)}: A geometric object classifying vector bundles on perfectoid spaces; central to the \emph{Fargues-Fontaine conjecture} (geometric Langlands for p-adic fields).
\item \textbf{Bhatt-Scholze (2019)}: Prismatic cohomology -- a unified cohomology theory for p-adic formal schemes.
\item \textbf{Scholze (2020+)}: Condensed mathematics -- extending perfectoid ideas to all of topology/algebra.
\end{itemize}

\textbf{Why This Matters Beyond Number Theory.} Prismatic cohomology, born from perfectoid spaces, is now the standard tool for $p$-adic Hodge theory and is being extended to singular and non-proper families. The condensed mathematics program is using perfectoid ideas to replace traditional topological spaces with a more flexible foundation, reaching into homotopy theory and quantum field theory far beyond the original number-theoretic setting.

\subsection{After the Proof}

The theory of perfectoid spaces is fully established. Scholze's construction of the tilt equivalence, the category of perfectoid spaces, and the pro-étale topology have all been rigorously developed and widely adopted. The foundational results---including the weight-monodromy conjecture (Bhatt-Morrow-Scholze, 2016) and prismatic cohomology (Bhatt-Scholze, 2019)---are complete.

Perfectoid spaces have transformed number theory by providing the bridge between mixed characteristic and equal characteristic $p$ that the field had long needed. Prismatic cohomology unified crystalline, de Rham, and étale cohomology into a single framework. Condensed mathematics, Scholze's later project, extended these ideas beyond number theory to all of topology and algebra, replacing traditional topological spaces with a more flexible foundation. The Fargues-Fontaine curve, built from perfectoid spaces, is now central to the p-adic Langlands program.

Major open questions remain. The geometric Langlands program for $p$-adic fields---the Fargues-Fontaine conjecture---is still wide open and would represent the culmination of $p$-adic Hodge theory in a geometric form. Understanding $p$-adic Hodge theory in mixed characteristic for singular or non-proper families continues to push the boundaries of the theory. The computational aspects of perfectoid spaces---explicit algorithms for computing with perfectoid rings and their cohomology---are in their infancy and could have significant applications to number-theoretic computations.

\section{Impact}
\begin{itemize}
\item \textbf{p-adic Langlands}: Perfectoid spaces provided the geometric framework for the p-adic local Langlands correspondence beyond $\text{GL}_2(\mathbb{Q}_p)$.
\item \textbf{Arithmetic geometry}: Prismatic cohomology unifies crystalline, de Rham, and \'etale cohomology.
\item \textbf{Arithmetic topology}: Condensed mathematics (Scholze, 2020) extends perfectoid ideas to all of homotopy theory.
\item \textbf{Computational}: Perfectoid methods give explicit algorithms for computing Galois representations.
\end{itemize}

\section{Exercises}

\begin{exercise}[Intermediate]
Explain the tilt equivalence: how does $K^\flat = \varprojlim_{x \mapsto x^p} K$ produce a field of characteristic $p$ from a field of characteristic $0$?
\end{exercise}

\begin{exercise}[Challenge]
Explain how the pro-\'etale topology solves the problem of doing p-adic Hodge theory in families. Why does the usual \'etale topology fail?
\end{exercise}

\section{Timeline}

\begin{tabularx}{\linewidth}{lX}
\toprule
\textbf{Year} & \textbf{Event} \\ \midrule
2011 & Scholze's PhD thesis introduces perfectoid spaces \\ 
2012 & Scholze proves local Langlands for $\text{GL}_n$ using Lubin-Tate perfectoid spaces \\

2013 & Scholze extends Fontaine's p-adic Hodge theory to rigid spaces \\
2016 & Bhatt-Morrow-Scholze develop $A_{\text{inf}}$-cohomology \\

2016 & Scholze-Weinstein: moduli of p-divisible groups as perfectoid spaces \\

2018 & Fargues-Fontaine curve; Fargues-Fontaine conjecture \\

2019 & Bhatt-Scholze: Prismatic cohomology \\

2018 & Scholze awarded Fields Medal \\

2020+ & Condensed mathematics extends perfectoid ideas to topology \\
2021 & Pila, Shankar, and Tsimerman prove the Andr\'e--Oort conjecture for Shimura varieties \\
2023 & Bakker, Brunebarbe, and Tsimerman prove the Griffiths conjecture on hypersurface cycles \\
2026 & Jacob Tsimerman awarded the Fields Medal for work on Shimura varieties and Hodge theory \\
\bottomrule
\end{tabularx}

\section{Citations}

\begin{enumerate}
\item Scholze, P. (2012). ``Perfectoid spaces.'' Publications Mathématiques de l'IHÉS, 116(1), 245--313.
\item Scholze, P. (2013). ``p-adic Hodge theory for rigid-analytic varieties.'' Forum of Mathematics, Pi, 1, e1.
\item Bhatt, B., Morrow, M., \& Scholze, P. (2016). ``Integral p-adic Hodge theory.'' Publications Mathématiques de l'IHÉS, 128(1), 219--397.
\item Scholze, P., \& Weinstein, J. (2016). ``Moduli of p-divisible groups.'' Cambridge Journal of Mathematics, 4(2), 145--211.
\item Fargues, L., \& Fontaine, J.-M. (2018). ``Courbes et fibrés vectoriels en théorie de Hodge p-adique.'' Astérisque, 406.
\item Bhatt, B., \& Scholze, P. (2019). ``Prisms and prismatic cohomology.'' arXiv:1905.08229.
\end{enumerate}
